\documentclass[11pt]{article}
\usepackage[a4paper,margin=27mm]{geometry}
\usepackage{amsmath,amssymb,mathtools}
\usepackage{fontspec,unicode-math}
\setmainfont{Latin Modern Roman}
\setmathfont{Latin Modern Math}
\usepackage{xurl,hyperref}
\hypersetup{colorlinks=true,urlcolor=blue,linkcolor=blue}
\newcommand{\sref}[2]{\hyperlink{source:#1:#2}{[S#2]}}
\newcommand{\cataloguescope}[1]{\paragraph{Mathematical question.} #1}
\begin{document}
\section{Conjecture 1.4 — Shorey}\label{statement}

\subsection{Definitions and mathematical statement}

\paragraph{Shared notation.}$\mathbb N=\{1,2,\ldots\}$, and $\mathbb Z,\mathbb Q,\mathbb R,\mathbb C$ denote the integers, rationals, reals and complex numbers. Any other conventions are specified locally.


Let $L\ge2$, $n_1>\cdots>n_{L-1}>0$, and $b_1,\ldots,b_L\in\mathbb Z\setminus\{0\}$, and write $f(X)=\sum_{i=1}^{L-1}b_iX^{n_i}+b_L$ with $H=\max_i|b_i|$. Assume that $f$ has at least two distinct complex roots. A proper vanishing subsum is a nonempty strict subset of the displayed signed terms whose sum equals zero.
\paragraph{Statement recorded in the supplied source.}
Is there $C=C(L,H)>0$ such that any integers $m\ge2$, $x$, and $y$ with $|y|>1$ and $y^m=f(x)$ satisfy either $m\le C$ or the expression $y^m-\sum_{i=1}^{L-1}b_i x^{n_i}-b_L$ has a proper vanishing subsum? \sref{statement}{2}
\subsection{Short English statement}
Is the perfect-power exponent bounded in terms only of the number and size of nonzero coefficients, unless cancellation occurs in a proper subset of the equation terms?
\subsection{Sources}
\begin{description}
\item[\hypertarget{source:statement:1}{S1}] ULAM.AI and the UnsolvedMath Contributors, UnsolvedMath: A Curated Collection of Open Mathematics Problems, catalogue entry. CC BY 4.0. \url{https://huggingface.co/datasets/ulamai/UnsolvedMath}.
\item[\hypertarget{source:statement:2}{S2}] Michel Waldschmidt, Open Diophantine Problems, arXiv:math/0312440v2 (24 January 2004), Conjecture 1.4 and immediately preceding sparse-polynomial hypotheses. \url{https://arxiv.org/abs/math/0312440}.
\end{description}
\label{end:statement}
\end{document}
